\section{Discussion}\label{sec:attack:discussion}

The results of the chapter are the following, each within the limits of its runs.
Fitted on paired sequences, an affine map attaches Pythia to the frozen hub: word-F1 rises from \vf{0.348} at \vf{10} paired sequences to \vf{0.876} at \vf{1{,}000}, against a hub reference of \vf{0.946} (H2; three seeds, Corpus A). % src: results/v7_2/main/mixed-v7_9-s4.json; factorial.json
On the unpaired pools alone, matching yields no readable map among those tested: no blind map exceeds its mismatched-prompt control by more than \vf{0.025} (H3).
Added to the supervised fit, the matching term raises word-F1 by \vf{0.309}, \vf{0.197}, and \vf{0.051} at \vf{10}, \vf{30}, and \vf{100} paired sequences, and by \vf{0.010} at \vf{300}, where the gain is still positive (H4, H5; three seeds, Corpus A). % src: results/v7_2/main/mixed-v7_9-s5.json
On legal text, a leak of \vf{90} rows yields a word-F1 of \vf{0.321} when spread over \vf{90} prompts, against \vf{0.178} when taken from two prompts, and no word class improves on a random draw of rows (H6; one seed, legal128w). % src: results/legal128w/v7_2/main/mixed-oct_w5-s1-L90.json; mixed-oct_w5b-L-prefix.json
H7 is contradicted: attached by its projector alone, a larger decoder lowers word-F1 on the same mapped states by \vf{0.273} against the best hub decoder (one seed). % src: results/legal128w/v7_2/main/oct_w5c-projector-llama3b.json

\paragraph{What the attacker reads.} The attack reads the frozen hub and its decoders, the hidden states of the target model for prompts without text, the attacker's own text of the same domain, and a leak of rows.
The leak is the only resource that ties a hidden state of the target model to text, and it requires that model's tokenizer for the leaked prompts.
With \vf{90} rows, about the content of ten short questions, the hybrid fit reconstructs Corpus A prompts at a word-F1 of \vf{0.60}; on legal text the same number of rows yields \vf{0.32}, with names and numbers substituted (Table~\ref{tab:attack:legal-examples}). % src: results/v7_2/main/mixed-oct_w5b-A-prefix.json (0.6042); results/legal128w/v7_2/main/mixed-oct_w5-s1-L90.json (0.3207)
The supervised fit on the same rows yields word-F1 scores of \vf{0.35} and \vf{0.11} on the two corpora; the difference comes from the unpaired pools, which carry no pairing between hidden states and text. % src: same files (0.3516, 0.1084)

\paragraph{Limitations.} Six limitations bound the claims.
\begin{itemize}
\item \textbf{Selection prompts.} The \vf{500} selection prompts chose the settings and report most scores.
Held-out scores exist only for leaks of one row per prompt on legal128w and for one row of Table~\ref{tab:attack:negative}; those prompts are held out from the attack settings, not from the hub checkpoint, which all validation prompts selected.
The test splits have not been read.
\item \textbf{Seeds.} The Corpus A budgets rest on three seeds; the leak study, the second target model, the larger decoder, and the rewriting model rest on one.
The intervals describe variation over prompts, not over seeds.
\item \textbf{Settings chosen with known text.} The supervision fraction, the temperature, and the batch size were selected by word-F1 on prompts whose text is known.
An attacker without such prompts would have to fix these settings in another way, and no run measures the cost.
The setting of the leak study is also not the best one found: a later grid on the selection prompts yields a setting that scores \vf{0.017} higher on the held-out prompts (interval \vf{0.013} to \vf{0.021}). % src: results/legal128w/v7_2/main/oct_w5-summary-confirmation-B.json decisions.P6/L90 (+0.0172 [0.0131, 0.0213])
\item \textbf{Same-domain pools.} Both unpaired pools and the leak come from the training prompts of the corpus under attack.
The results do not show how the attack behaves when the attacker's text comes from another domain, or when the unpaired pool is small.
\item \textbf{Models and layer.} The results cover one block of two target models of about one billion parameters, and three hub models.
They do not cover a deployed system, a larger model, a quantized or noised hidden state, or a defense.
\item \textbf{Metric.} Word-F1 measures word overlap and ignores order.
It does not separate a correct reconstruction from a fluent one with substituted names and numbers, and no measure of meaning is reported.
\end{itemize}

The leak supplies correspondences that constrain the affine map, and the matching term compares unpaired representations; the experiments do not establish that either source determines the map alone.
\todo{Author: the information-theoretic framing that Chapter~\ref{ch:future} refers to is not written.}
